% spellchecker:disable
\documentclass[pdftex,
	chapterprefix,
	headsepline,
	footsepline,
	% colordvi,
	11pt,
	a4paper,
	parskip=half,
	% enabledeprecatedfontcommands,
	final,
	appendixprefix,
	bibliography=totoc]{scrbook}

% define a new if command to distinguish between PDF and DVI output
\usepackage{ifpdf}
\ifx\pdfoutput\undefined
\pdffalse %not PDFLaTeX
\else
\pdfoutput=1
\pdftrue
\fi

% language support (hyphenation etc)
% this sets the main language of the paper to english
% if you wish to compose it in german, switch the options around
\usepackage[german,english]{babel}

% for prettier tables
\usepackage{booktabs}

\usepackage[section]{placeins}

% support for latin1 characters. That means you can enter umlauts directly
% no need for "a "u "o "s anymore
\usepackage[utf8]{inputenc}
\usepackage{textcomp}

% provides the \url{} command to pretty print urls
\usepackage{url}

% allows to \includegraphics
\usepackage{graphicx}
\usepackage{svg}

% defines some standard colornames like "black" etc.
\usepackage{color}

% allows to color tablecells
\usepackage{colortbl}

% provides an easier interface to if-then-else constructs in custom macros
\usepackage{ifthen}

% allows tables to break over pages.
\usepackage{supertabular}

% allows to have different kinds paper orientations in the same pdf-documnent
\usepackage{pdflscape}

% allows to specify absolute texpos for textboxes. This is generally only important for the titlepage
\usepackage[absolute]{textpos}

% % allows to enumerate different figures with a) b) in the same figure-environment.
\usepackage{subfigure}

\usepackage{caption}
\DeclareCaptionLabelFormat{cont}{#1~#2\alph{ContinuedFloat}}
\captionsetup[ContinuedFloat]{labelformat=cont}

% math
\usepackage{amsmath}

% allows inserting source code
\usepackage{listings}
\usepackage{scrhack} % fixes lst incompatibility
\lstset{
% numbers=left,              % Location of line numbers
	numberstyle=\tiny,          % Style of line numbers
	numbersep=5pt,              % Margin between line numbers and text
	tabsize=4,                  % Size of tabs
	extendedchars=true,
	breaklines=true,            % Lines will be wrapped
	showspaces=false,
	showtabs=false,
	showstringspaces=false
}

% More fancy syntax highlighting: https://texdoc.net/texmf-dist/doc/latex/minted/minted.pdf
% careful, this has additional external dependencies
%\usepackage[newfloat]{minted}
% configure source code highlighting; edit this to your liking
% style=bw will save you some money on color pages, but may be harder to read
%\setminted{
%	style=stata-light,
%	frame=leftline,
%	breaklines=true,
%	numbers=left,
%	tabsize=4}

% allows usage of BibLaTeX for bibliography; needs the Biber tool!
\usepackage[
backend=biber,
sorting=anyt
]{biblatex}
\usepackage{csquotes}
% point this to your bibliography file
\addbibresource{bib/references.bib}

% finetune the gaps between figure and text in the subfigure environment (basically close the gap as much as possible)
\renewcommand{\subfigtopskip}{0pt}
\renewcommand{\subfigbottomskip}{0pt}

% some color definitions for the pdf statements below
\definecolor{mygrey}{rgb}{0.45,0.45,0.45}
\definecolor{mydarkgrey}{rgb}{0.2,0.2,0.2}
\definecolor{red}{rgb}{1.0,0.33,0.33}
\definecolor{orange}{rgb}{1.00,0.73,0.33}
\definecolor{yellow}{rgb}{0.95,0.92,0.}
\definecolor{lightgreen}{rgb}{0.3,0.95,0.46}
\definecolor{titleblue}{rgb}{0.03,0.10,0.46}

\ifpdf
% For screen viewing it is nice to have references marked in a slightly different
% color than the rest of the text. Since they will be hyperlinks to the
% referenced objects.
\usepackage[pdftex,
	pdftitle={Fancy thesis template},
	% colorlinks,
	% linkcolor={black},
	% citecolor={black},
	urlcolor={black},
	plainpages={false},
	bookmarksnumbered={true},
	pdfauthor={Vorname Nachname},
	pdfsubject={},
	pdfkeywords={},
	pdfstartview={FitBH}]{hyperref}

\pdfcompresslevel=9
\fi

% some configuration for the amount of text on a single page
\usepackage{typearea}
\areaset[1.5cm]{418pt}{658pt}
\setlength{\headheight}{37pt}

% enter author and title for the titlepage.
\author{}
\title{}

% To avoid nasty mistakes like having comments directly in the textflow
% the following \todo macro was defined. With that you can enter
% \todo{What I still have to do here}
% inside of your text and a marker will appear at the page's margin with the
% text "What I still have to do here".
% The first line activates this feature. If you comment it out and uncomment
% the second line below there will be no error messages and no todos will be shown
% anymore. So - even if you have forgotten to delete one of them - they will not appear
% in the final printout.
\newcommand{\todo}[1]{\marginpar{\textcolor{red}{TODO:} #1}}

% We recommend to split your document into several files. Usually one for every chapter is a
% good idea. If you follow this guideline (how to assemble these files in a single document
% see two paragraphs below) you will be able to single out chapters via the \includeonly{}
% command. Using this mechanism page numbering and references of the full run before will be
% preserved. This also nice, if your latex run tends to get slow and you need to fine tune
% some formatting in one chapter - just include that one. The rest (or at least the ones before
% the one currently under construction) will remain untouched. This means a boost in compilation time.
%\includeonly{content/chapter2}

\begin{document}

% the next two lines influence the detailedness of the table of contents
% and to what structure depth numbers are written before sections/subsections/paragraphs
% You should not touch this
\setcounter{tocdepth}{2}
\setcounter{secnumdepth}{3}

\frontmatter

% here the titlepage is included. Look into the file "frontpage.tex" to
% adapt it to your needs (name, title etc.)
% spellchecker:disable
\begin{titlepage}
\selectlanguage{german}
\vspace*{-1cm}
\newlength{\links}
\setlength{\links}{0.9cm}
\setlength{\TPHorizModule}{1cm}
\setlength{\TPVertModule}{1cm}
\textblockorigin{0pt}{0pt}

\sffamily
\LARGE

\begin{textblock}{16.5}(2.8,2.6)
 \hspace*{-0.25cm} \textbf{UNIVERSITÄT DUISBURG-ESSEN} \\
 \hspace*{-1.15cm} \rule{5mm}{5mm} \hspace*{0.05cm} FAKULTÄT FÜR INGENIEURWISSENSCHAFTEN\\
 \large{}ABTEILUNG INFORMATIK UND ANGEWANDTE KOGNITIONSWISSENSCHAFT\\
\end{textblock}


%Hier Titel, Name, und Matrikelnummer eintragen, \\ make a newline
\begin{textblock}{14.5}(3.2,8.5)
  \large
{ \textbf{Bachelorarbeit}} \\[1cm]
{\LARGE \Large\textbf{An Algorithm for Dependency-Preserving Smart Home Updates}} \\[1.3cm]
Koray Uzun\\
Matrikelnummer:3081690\\
Angewandte Informatik (Bachelor)
\end{textblock}



\begin{textblock}{10}(10.5,17.5)
\includegraphics[scale=1.0]{images/unilogo.pdf}\\
\normalsize
\raggedleft
Fachgebiet Verteilte Systeme, Abteilung Informatik \\
Fakultät für Ingenieurwissenschaften \\
Universität Duisburg-Essen \\[2ex]

\today\\[15ex]
\raggedright
% Supervisors
{\textbf{Erstgutachter:}} Prof. Dr-Ing. Torben Weis \\
{\textbf{Zweitgutachter:}} Peter Zdankin \\
{\textbf{Zeitraum:}} 1. September 2042 - 1. Januar 2043 \\
\end{textblock}

\end{titlepage}


\selectlanguage{english}

\chapter*{Abstract}


Smart Home Systeme sind Netzwerke von elektronischen Geräten innerhalb eines Wohnraumes. 
Diese Systeme unterliegen, wie jede andere Technologie auch, einem ständigen Wandel. Während andere Technologien,
wie zum Beispiel Smartphones meist nach einigen wenigen Jahren ausgetauscht, erwartet man von
Smart Home Systemen eine lange Lebensdauer.
Die meist komplexen und teuren Installationen sind nicht dazu ausgelegt alle paar Jahre 
ausgetauscht zu werden. Dennoch erwartet man, wie von anderen Technologien auch, dass die Systeme
sich den Gegebenheiten der Zeit anpassen. Neue Funktionen oder oft auch Sicherheitslücken erfordern 
ein ständiges updaten der Systeme oder einzelner Geräte. Mit diesen Updates geht jedoch ein großes Risiko 
einher. Aufgrund der riesen Diversität verschiedener Smart Home Systeme kann vom Hersteller meist nicht
garantiert werden, dass ein Update keine Schäden anrichtet. Lediglich das Update im Einzelnen wird auf 
mögliche Fehlerquellen überprüft, nicht jedoch das Zusammenspiel des Updates mit anderen Geräten.
Dies wäre aufgrund der riesen Diversität an Smart Home Systemen für den Hersteller auch gar nicht umsetzbar.
Dennoch benötig man eine Möglichkeit Updates präventiv zu überprüfen, um so größere Schäden zu meiden und 
die Langlebigkeit von Smart Home Systemen zu gewährleisten.
Diese Abschlussarbeit beschäftigt sich mit der Implementation und Evaluation eines Algrorithmus, welcher 
Abhängigkeiten zwischen Geräten untersucht und basierend auf diesen Abhängigkeiten Updatekonfigurationen 
findet, die keine Probleme verursachen.
Der Algorithmus basiert auf dem Paper "An Algorithm for Dependency-Preserving 
Smart Home Updates" von Peter Zdanking, Matthias Schaffeld, Marian Waltereit, Oskar Carl und Torben Weis. 






\footnote{Wikipedia: \url{https://en.wikipedia.org/wiki/Abstract_(summary)}}



\tableofcontents

%\listoffigures
\mainmatter

% To assemble the whole document
% Please be aware that each file will begin on a new page
% therefore chapters should be put into such a file.
% There cannot be an include statement inside of an "included" file.
% So if you want to further divide your document - use \input inside of
% the included files. \input will not begin on a new page.
\chapter{Einführung}

In den letzten Jahren stieg die Anzahl der Smart Home Geräte stark. Laut Verified Market Research wird
der gloable Smart Home Markt im Jahr 2019 auf bereits 80 Milliarden Dollar geschätzt. Prognostiziert werden 207.88 
Milliarden Dollar im Jahr 2027. Dies entspricht einer jährlichen Wachstumsrate von 13.52 Prozent [2].
Zudem ist es so, dass innerhalb von Smart Home Systemen die Anzahl der Geräte ebenfalls stetig steigt. Während 2016
durchschnittlich 5.8 Geräte über ein Smart Home System vernetzt waren, waren es 2018 schon bereits 8.1 Geräte [3]. 
Es ist zu vermuten, dass sich dieser Trend vorerst so weiterentwickeln wird und somit Smart Home Systeme immer
größer werden. Wie jede andere Technologie müssen Smart Home Systeme regelmäßig geupdatet werden, um zum Beispiel
Sicherheitslücken zu schließen oder neue Funktionalitäten hinzuzufügen. 
Mit der steigenden Anzahl an Geräten, steigt in einem Smart Home System auch logischerweise die Anzahl der Updates und somit 
insgesamt die Komplexität von Smart Home Systemen. Das Problem hierbei ist, dass mit steigender Komplexität zum einen automatisch eine höhere
Fehleranfälligkeit bezüglich Updates einhergeht, zum anderen wird es schwieriger Fehler zu beheben,
die durch Updates verursacht werden. Updates im Einzelnen besitzen bereits viele Fehlerquellen, wie zum Beispiel Sicherheitslücken 
oder generell fehlerhafte Implementationen. Dies sind jedoch Aspekte die bereits vom Entwickler weitestgehend vermieden werden.
Viel Problematischer ist das Zusammspiel verschiedener Geräte und den dazugehörigen Updates. Ein an sich fehlerfreies Update
kann bei der Installtion in einem Smart Home System, trotz "fehlerfreier" Programmierung große Probleme verursachen. Was ist wenn zum Beispiel 
ein Update eine Funktion eines Geräts so verändert, dass ein anderes Gerät nicht mehr darauf zugreifen kann. Durch solche kleine Änderungen
können Abhängigkeiten zwischen Geräten zerstört werden, wodurch ganze Teile eines Smart Home Systems ausfallen
können. Solch ein Ausfall bedeutet für den Nutzer meist ein großer finanzieller Schaden, da Updates oft nicht 
rückgängig gemacht werden können und somit Ersatz beschafft werden muss.
Aus diesen Gründen sind präventive Maßnahmen notwendig,  um die Langlebigkeit von Smart Home Systemen gewährleisten zu können.
Langlebigkeit wird nämlich von den meisten Usern erwartet, wie man an der Resonanz zum Abschalten der HueBridge V1 von Philips sehen kann 
[4].


\section{Ziel der Arbeit}


Eine Steigerung der Lebensqualität ist das Kernziel von Smart Home Systemen. Dies erreicht man zum Beispiel durch die Konfiguration von 
automatisierten Prozessen. Dadurch können Nutzer nämlich Zeit, Energie und in gewissen Belangen auch Geld sparen. Laut Umfragen geben 57 Prozent der Nutzer von
Smart Home Systemen in Amerika an jeden Tag durchschnittlich 30 Minuten Zeit einzusparen [4].
Zusätzlich lassen sich nach Berechnungen des Fraunhofer Instituts für Bauphysik  durch Automatisierungen bis zu ca. 40 Prozent Heizkosten 
sparen. [5].
Solche Vorteile zeigen, dass sich die anfangs teure Anschaffung eines Smart Home Systemes auf lange Zeit sogar rentieren kann. Sicherlich 
ist dies der Idealfall, aber dennoch zeigen diese Zahlen, dass Smart Home Systeme viele Potenziale bieten. Es hängt jedoch von 
vielen Faktoren ab, ob sich die Anschaffung eines Smart Homes lohnt und objektiv ist dies nicht zu beurteilen. Was jedoch 
objektiv gesagt werden kann, ist dass in den meisten Fällen Nutzer aufgrund der teuren Anschaffungskosten eines Smart Homes davon ausgehen, dass 
ihr System viele Jahre erhalten bleibt. Zum Beispiel gibt es Smart Home Systeme, welche beim Bau eines Hauses direkt mit geplant werden,
sodass man, wie zum Beispiel bei KNX Systemen alle Geräte über ein permanent installiertes Bus System verbinden kann [1].
Selbstverständlich gibt es auch günstigere Varianten, aber dennoch ist die generelle Erwartungshaltung, dass wenn man ein funktionierendes System hat, dieses
auch über viele Jahre hinweg weiter funktionieren wird.
Daher is das Ziel dieser Arbeit die Implementation und Evaluation eines Algrorithmus, welcher 
Abhängigkeiten zwischen Geräten untersucht und basierend auf diesen Abhängigkeiten Updatekonfigurationen 
ermittelt, welchen den größten Nutzen für den Nutzer haben.
Mit diesem Algorithmus wäre es für Smart Home Nutzer möglich Updates bereits vor der Installation zu überprüfen, um so sicherzustellen, 
dass sie keine Ausfälle verursachen. Die Anwendung dieses Algorithmus sollte für den Nutzer möglichst einfach sein. So könnte man zum 
Beispiel ein Smartphone nutzen, welches als zentrales Gerät im System dient. Auf diesem zentralen Gerät wäre es dann möglich den 
Algorithmus laufen zu lassen und sich über das Smartphone dann für die passende Updatekonfiguration zu entscheiden. Smartphones
bieten sich besonders an, da sie eine intuitive Bedienung für den Nutzer ermöglichen und im Normalfall jeder Nutzer eines Smart Homes
ein Smartphone besitzt. So umgeht man das Problem, dass es, wie bereits erwähnt, etliche Smart Home Anbieter mit verschiedener
Software und Geräten gibt. 




 


\chapter{Related Work}\label{ch:related_work}
In diesem Kapitel wird das Thema Langlebigkeit im Bereich Software Engineering thematisiert. Infolgedessen werden
Smart Home Systeme unter dem Aspekt der Langlebigkeit genauer beleuchtet und  das Thema Dependency Management aufgegriffen, wobei bereits 
bestehende Ansätze/Lösungen diesbezüglich genauer betrachtet werden. Zusätzlich wird der aktuelle Stand des Dependency Managements im 
Bereich Smart Home untersucht.


\section{Smart Homes}
Smart Home Systeme sind grundlegend nichts anderes als die Verknüpfung elektronischer Geräte. Durch die hohe Anzahl 
der Anbieter für Smart Home Systeme existieren jedoch viele Architekturen, die teilweise sehr unterschiedlich funktionieren .
Im Rahmen dieser Arbeit ist es nicht notwendig die verschiedenen Smart Home Architekturen genauer zu beleuchten,
dennoch schadet es nicht sich einen groben Überblick über die aktuellen Technologien zu verschaffen. Grob können die
verschiedenen Architekturen darin unterscheiden werden, ob sie eine Verknüpfung zu einer Cloud besitzen. Cloud gebundene
Architekturen bieten dem Nutzer meist eine einfachere Installation und Wartung, indem zum 
Beispiel Updates automatisch installiert werden. Gleichzeitig bringen sie aber auch potenzielle Gefahren mit sich, da die
Abhängigkeit zu einer Cloud Smart Home Systeme verwundbar machen kann. Wird ein externer Service, wie zum
Beispiel die Cloud abgeschaltet, führt dies bei Cloud-gebundenen Archtiekturen oft zu Ausfällen im System[7].
Des Weiteren sind Geräte, die mit einer Cloud verbunden sind anfälliger gegenüber Angriffen von Außenstehenden, die sich 
Zugriff auf das System beschaffen wollen.
Nicht Cloud-gebundene Architekturen haben diese Nachteile nicht. Es ist jedoch meist schwieriger eine solche Architektur aufzubauen, 
da zum Beispiel, wie im Falle von OpenHab[?] viel manuelle Konfiguration benötigt wird. Zusätzlich bieten nicht Cloud-gebundene Architekturen
mehr Sicherheit, da sie keine starken Abhängigkeiten zu anderen Dienstleistern benötigen.
Natürlich könnte man Smart Home Systeme noch weiter spezifizieren, indem man einzelne Komponenten genauer beleuchet und 
verschiedene Smart Home Anbieter miteinander vergleicht, jedoch ist dies, wie bereits erwähnt, im Rahmen dieser 
Arbeit nicht notwendig. Der zu entwickelnde Algorithmus soll später theoretisch Platform- und Architekturunabhängig anwendbar sein.
Zu diesem Zweck werden im Folgenden einige Aspekte bezüglich Smart Homes definiert und vorausgesetzt. 
Zdankin beschreibt in seiner Arbeit Smart Home Systeme als:

\par
\begingroup
\leftskip=1cm % Parameter anpassen
\noindent %ab hier der Text, der eingerückt werden soll
\textit{"A smart home has a set of devices that are connected through
a platform. Each device has a certain software version and
a set of available updates. Each software version has a set
of available predefined services. Devices can use services of
other devices which creates dependencies. An update configuration
is one of the finite states of the nondeterministic
finite automaton (NFA) that can be constructed by using the
configurations as states and connecting them using individual
updates as transitions." (Zdankin et al., 2020, Requirements and Mechanisms for Smart Home Updates [8])}
\par
\endgroup
Um Missverständnisse zu vermeiden folgen nun zusätzlichen Begriffserklärungen.
Die folgenden Elemente in Abbildung 2.1 sind für ein Smart Home System also essentiell:
\begin{itemize}
\item Device (Gerät): Geräte sind die Komponenten, die die eigentliche Funktionalität anbieten. Sie besitzen eine aktuelle Firmware
Version mit entsprechenden Diensleistungen(Services) und potenzielle Updates. Zusätzlich können sie Abhängigkeiten zu anderen Geräten haben, indem sie auf deren
Dienstleistungen zugreifen.
\item Firmware Version: Die Firmware Version gibt an, auf welcher Version ein Gerät sich befindet und gibt somit an welche Dienstleistungen ein Gerät anbietet.
\item Service (Dienstleistungen): Eine Dienstleistung ist eine Funktionalität, die ein Gerät anbietet. Zum Beispiel bietet ein Thermostat
die Dienstleistung "Temperatur messen" an.
\item Update: Eine Aktualisierung der Firmware eines Geräts. Durch ein Update können sich angebotenen Diensleistungen eines Geräts ändern.
\item Platform: Die Komponente, die alle Geräte miteinander vernetzt.
\item Dependency (Abhängigkeiten): Eine Abhängigkeit zwischen zwei Geräten besteht, wenn ein Gerät Zugriff auf die Dienstleistungen eines
anderen Geräts benötigt.
\end{itemize}

\begin{figure}[!ht]
\hspace*{-2,5cm}                                                           
\includegraphics{"System Model"}
  \caption{Grundlegende Komponenten eines Smart Home Systems}
  \label{fig:Prob1:MEA}
\end{figure}


\newpage
\subsection{Prescriptive vs Descriptive Standards}
Für die Kommunikation der Geräte untereinander existieren viele Kommunikationsprotokolle wie zum Beispiel Bluetooth,
Wi-Fi oder auch XMPP. Diese sind bereits im großen Umfang in Verwendung und stark erforscht, weshalb es keiner weiteren
Untersuchung dieser Protokolle im Rahmend dieser Arbeit bedarf.
Kommunikation allein recht jedoch nicht, die Geräte müssen nicht nur untereinander kommunizieren, sondern sie müssen sich auch 
verstehen. Dafür ist ein Standard notwendig, welcher die Funktionen/Dienstleistungen der Geräte beschreibt.
Die zwei führenden Ansätze im Bereich IoT/Smart Homes sind deskriptive und preskriptive Standards.
Zdankin beschreibt die beiden Standards in seiner Arbeits folgendermaßen:

\textit{
"A standard for service definitions that only regulates how a device can describe itself is
a descriptive standard, because it allows vendors to describe their devices in their own
terms. A recently finalized descriptive standard is the Web Of Things [Ka20; Ko20]. If
devices need to communicate across different descriptive standards, a translation must be
considered. However, McCool et al. have stated security concerns about purely descriptive
approaches, such as scanning for door locks with known physical weaknesses [MR18] (Zdankin et al., 2020, Requirements and Mechanisms for Smart Home Updates [8]).}
\todo{Wie muss ich die Quellenangaben im Zitat zitieren?}

%\footnotetext{[Ka20]: httpsȷ//www.w3.org/TR/wot-thing-description/}
%\footnotetext{[Ko20]: httpsȷ//www.w3.org/TR/wot-architecture/}
%\footnotetext{[MR18]: Mccool, M.; Reshetova, E.ȷ Distributed Security Risks and Opportunities in the
%W3C Web of Things. In. Jan. 2018.}

\textit{
A prescriptive standard does not only regulate the how but also the what of self-description.
These standards prescribe how devices and services should be defined and vendors can use
these terms for unambiguity. By using the same terms to define services, translation is not
required anymore. A prescriptive standard defines all possible devices and all their services
they could implement. However, a device can implement a subset of these services only.
This restriction has benefits as well because it ensures compatibility across device vendors.
Widely used examples are the smart home standards of Amazon, Apple, and Samsung6."
(Zdankin et al., 2020, Requirements and Mechanisms for Smart Home Updates [8])}

In der Realität bieten sich hybride Standards vermutlich am meisten an, da sowohl deskriptive als auch preskriptive Standards Vor- und
Nachteile besitzen.   


\subsection{Update Planing}
Oft werden Smart Home Geräte automatisch über eine aktive Internetverbindung geupdatet oder auch durch andere
Geräte wie zum Beispiel einem Hub. Dies bietet dem Nutzer mehr Komfort, da kein manueller Aufwand für Updates ensteht.
Solches "Remote Updating" sollte aber eigentlich vermieden werden, da dadurch zum Beispiel schädliche Software installiert werden
kann. Außerdem müsste der externe Diensleister, der das Updaten der Geräte übernimmt, viele Informationen über das System sammeln,
um so ideale Updatekonfigurationen zu finden. Das Preisgeben aller Informationen sollte nicht leichtsinnig getan werden, da es viele
potenzielle Gefahren birgt. Allein die Information, wann eine automatisierte Heizung anfängt zu heizen kann potenziellen Kriminellen Informationen
darüber geben, wann jemand zu Hause ist oder wann ein Haus leer steht. 
Idealerweise sollten Nutzer daher das Updaten lokal selbst übernhemen. Dies bedeutet selbstverständlich mehr manuelle Konfiguration,
vor jedem Update, da der Nutzer gefragt werden sollte, ob er ein Update installieren möchte. Diese Frage kann ein Nutzer ohne weitere Informationen
gar nicht plausibel entscheiden. Es wäre also von Vorteil, wenn es für den Nutzer
eine Übersicht über alle möglichen Updates geben würde, in der er sich dann für passende Updates entscheiden kann. 
Für einen solchen Zweck bietet sich Smartphone als zentrales Gerät an um eine Übersicht über alle im 
System vorkommenden Geräte zu erstellen.
Dazu benötigt man eine Möglichkeit die Informationen über die Services und Updates eines Geräts leicht abfragen zu können.
Eine solche Abfrage müsste in der Realität für jedes Gerät individuell angepasst werden. Dies ist jedoch nicht umsetzbar,
weswegen es nötig ist einen einheitlichen Weg zu finden.
Zdankin schlägt vor die Metadaten der Geräte zu diesem Zweck zu nutzen. Geräte besitzen üblicherweise Metadaten, welche Informationen über 
Software Version, Speicherbedarf usw. enthalten. Diesen Metadaten könnte man zusätzlich Informationen über Dienstleistungen und Updates 
hinzufügen, sodass man einen schnellen Überblick über alle Geräte innerhalb eines Netzwerks erhalten kann
(Zdankin et al., 2020, Requirements and Mechanisms for Smart Home Updates [8]). 

\section{Langlebigkeit}
Im Bereich Software Engineering ist Langlebigkeit ein bekanntes Problem. Es wurden bereits Untersuchungen über die durchschnittliche 
Lebensdauer von Software durchgeführt mit dem Ergebnis, dass kleine Projekte im Bereich von unter 100.000 Zeilen Code wird eine Lebensdauer von circa 
6-8 Jahren haben. Dabei ist zu beobachten, dass mit steigender Größe eines Projekts die erwartete Lebensdauer im Durchschnitt
ebenfalls steigt [9]. Dies hängt vermutlich mit den hohen Kosten zusammen, die für größere Projekte aufkommen, weswegen man daran
interessiert ist lange einen Nutzen von der Software zu haben. Gründe für das Altern von Software sind vielseitig. So kann zum Beispiel 
die Wartung von Software einfach zu teuer werden, so dass es sich lohnt neue Software zu entwickeln oder eine Entwicklungsumgebung wird
nicht mehr unterstützt und es lohnt sich nicht die Software zu portieren. Andere Gründe können Sicherheitslücken oder auch fehlerhafte 
Updates sein. Die German Research Foundation hat ein Projekt namens Design for Future ins Leben gerufen, um das Problem der Langlebigkeit
von Software zu untersuchen. Das Projekt beschreibt sein Ziel darin fundamentale neue Ansätze im Bereich langlebiger Software zu finden, um Probleme mit Legacy-Software
oder der Adaption von Software auf neue Plattformen zu lösen und eine kontinuierliche Weiterentwicklung von Software Systemen gewährleisten 
zu können [10].
Im Bereich Smart Home existieren keine solche Programme. Das Problem der Langlebigkeit wird in diesem Kontext unterschätzt oder einfach ignoriert.
Aufgrund der Komplexität der Systeme besitzen sie viele Schwachstellen die zu einem frühzeitigen Ausfall des Systems führen können. Dazu gehören:

\begin{figure}[h]
\centering
\includegraphics{"Threats"}
\caption{Gefahren für die Langlebigkeit von Smart Home Systemen\( ^1\) }
\label{fig:Prob1:MEA}
\end{figure}
\footnote{Zdankin et al, 2020, An Algorithm for Dependency-Preserving Smart Home Updates}

In dieser Arbeit wird hauptsächlich die Gefahr der "Breaking Updates" thematisiert.
"Breaking Updates" können Abhängigkeiten im System zerstören und sind somit Schäden gegen die man als Nutzer nichts unternehmen
kann. Eine Möglichkeit wäre es Updates generell nicht zu installieren, wodurch Nutzer auf neue Funktionalitäten verzichten müssten, während gleichzeitig das Risiko steigt,
dass Sicherheitslücken im System aufkommen. Updates also grundlegend zu ignorieren ist nicht empfehlenswert, weswegen es einer Möglichkeit bedarf
Updates vor ihrer Installation zu überprüfen.
Eine andere Gefahr ist zum Beispiel, dass ein externer Service, der für das Smart Home System notwendig ist, offline gehen kann (Discontinued External Service).[1]
Der Unteschied zu den "Breaking Updates" besteht darin, dass man als Käufer eines Cloud-gebundenen Smart Home Systems diese Gefahr aktiv und willentlich eingeht.
Die ist bei "Breaking Updates" nicht der Fall.
Außerdem lässt sich das Ausmaß der Gefahr eines "discontinued Services" weitesgehend vermeiden, indem man bei etablierten Anbietern einkauft oder
komplett auf die Abhängigkeit zu einem externen Service verzichtet. Auch dies ist bei "Breaking Updates" nicht möglich.


\newpage
\section{Dependency Managment}
Unter Dependency Management versteht man im Allgemeinen das Strukturieren von Abhängigkeiten verschiedener 
Systeme/Programme. Oft können Programme nur in Abhängigkeit von anderen Programmen starten.
Dies stellt kein Problem dar, solange die Abhängigkeiten in ihrer Anzahl überschaubar bleiben. 
Realität ist jedoch, dass Programme meist von mehreren Programmen abhängig sind und diese
Prorgramme wiederum abhängig von anderen Programmen. Bei zu vielen Abhängigkeiten verliert man 
als Entwickler und auch als Nutzer schnell den Überblick. Dieses Wirrwarr an Abhängigkeiten bezeichnet man umgangssprachlich 
als "Dependency Hell", denn sobald man einmal den Überblick verloren hat, wird es äußerst schwierig weitere Änderungen am
System vorzunehmen.

\subsection{Probleme}
Es folgt eine kleine Zusammenfassung, der am häufigsten auftretenden Formen der "Dependency Hell" [11]:
\begin{itemize}
 \item Long Chain of Dependencies: Ein Programm A ist abhängig von einem Programm B. Das Programm B wiederum ist abhängig von einem Programm C.
Möchte man nun Programm  A nutzen benötigt man zusätzlich Programm  B und damit auch Programm C. Solche lange Ketten können Konflikte verursachen.
(Siehe Conflicting Dependencies)
\item Conflicting Dependencies: Sei ein Programm A abhängig von einem Programm B1.1, ein anderes Programm B ist abhängig von  Programm B1.2. Programm B 
kann aber nicht gleichzeitig mit der Version 1.1 und der Version 1.2 installiert sein. Dies bedeutet, Programm A und B können nicht gleichzeitig in einem System laufen.
\item Circular Dependencies: Sei ein Programm A abhängig von einer spezifischen Version von Programm B. Programm B wiederum ist abhängig von einer spezifischen
Version von Programm A. Beide Programme können nun nicht geupdatet werden, da sie die Abhängigkeiten kaputt verletzen würden.
\item Version Lock: Durch einen Version Lock ist man dazu gewungen beim updaten einer Software alle abhängigen Programme ebenfalls upzudaten. (semver.org)
\item Version Promiscuity: Werden Abhängigkeiten zu locker betrachtet können zukünftige Updates als kompatibel angesehen werden, obwohl sie es nicht sind.
\end{itemize}

Für die genannten Probleme existieren bereits viele Ansätze, um gegen sie vorzugehen oder sie zu vermeiden. Dazu gehören [11]:

\subsection{Lösungen}
\begin{itemize}
 \item Semantic Versioning [12]: Semantic Versioning versioniert Software nach dem Format "MAJOR.MINOR.PATCH", um Probleme wie Version Lock 
oder Version Promiscuity zu vermeiden. Dafür benötigt 
man eine Schnittstelle der man die Änderungen an der Software mitteilt. Bei Änderungen, die die API betreffen und nicht abwärtskompatibel
sind muss die MAJOR Nummer erhöht werden. Bei Änderungen, die abwärtskompatibel sind, muss die MINOR Nummer erhöht werden. 
Änderungen, die die API nicht betreffen erhöhen die Patch Nummer. Es gibt noch viele weitere Spezifikationen bezüglich Semantic
Versioning, das grobe Konzept sollte jedoch verstanden sein. 
\item Package Manager: Package Manager sind Tools die sich um das installieren, updaten und konfigurieren von Programmen automatisch
kümmern, wodurch eine manuelle Konfiguration nicht mehr nötig ist.
\item Portable Applications: Portable Applikationen umgehen die Probleme, die
Abhängigkeiten mit sich bringen. Sie funktionieren selbstständig,
indem alle nötigen Komponenten bereits vorhanden sind.
\end{itemize}

Die Auflistungen sollen zeigen, dass Dependency Management ein allgegenwärtiges Thema ist, welches in den meisten Bereichen bereits
viel Aufmerksamkeit bekommt. Dependency Management umfasst natürlich noch viele weitere Probleme und dazugehörige Lösungen und Lösungsansätze,
im Bereich Smart Home exisitiert jedoch nahezu keine Forschungsgrundlage. Als Nutzer kann man nur hoffen, dass neue Updates keine Schäden anrichten, 
da es keine Möglichkeit gibt sich davor zu schützen. 


\section{Fazit}

Im Bereich Software Engineering wird Langlebigkeit viel Aufmerksamkeit geschenkt. Smart Home Systeme sind selbstverständlich Teil des
Software Engineerings, aber schaut man sich den aktuellen Stand der Forschung der Langlebigkeit explizit im Bereich der Smart Home
Systeme an, wird deutlich das diesbezüglich kaum Forschungsgrundlagen bestehen.
Dabei werden Smart Home Systeme als Luxusgüter angesehen von denen man aufgrund der hohen Preise eine lange Lebensspanne erwartet.
Um dies zu gewährleisten werden generelle Ansätze des Software Engineerings bezüglich Langlebigkeit genutzt, ausreichend ist dies jedoch nicht.
Ohne vernünftiges Dependency Management kann nicht gewährleistet werden, dass Updates keine Ausfälle in Smart Home Systemen verursachen.
Dies sieht man am Beispiel des Logitech Harmony Hub's [13]. Logitech hatte eine \todo{umschreiben}
neue Firmware Version für das Hub veröffentlicht, welche Sicherheitslücken schließen und Bugs fixen sollte.
Nachdem Update kam es vermehrt zu Ausfällen zwischen dem Hub und third-party Geräten. Den Nutzern waren die Hände 
gebunden, sie konnten nur darauf hoffen, dass Logitech die Probleme wieder behebt. Dies hat Logitech in diesem Fall auch getan,
dennoch zeigt dieses kleine Beispiel, wie leichtsinnig Updates von Nutzern installiert werden. Eine andere Möglichkeit besitzt man auch als 
Nutzer eines Smart Homes nicht wirklich, da es wie eingangs erwähnt, zur Zeit keine Möglichkeiten gibt Updates vor der Installation zu 
überprüfen.








\chapter{Design und Implementation}\label{ch:eval}

Dieses Kapitel umfasst eine detaillierte Beschreibung des Designs und der Implementation des Algorithmus. Zunächst wird der grobe
Ablauf des Algorithmus beschrieben, um spätere Design- und Implementationsentscheidungen besser nachvollziehen zu können.
Zusätzlich wird ein Generator zum Erstellen von Smart Home Systemen vorgestellt und der mitgegebene Code beschrieben.


\section{Ablauf des Algorithmus}
Zunächst folgt eine grobe Beschreibung aller Schritte des Algorithmus, um so als Leser genauere Entscheidungen im Laufe der Arbeit
besser nachvollziehen zu können. Der Algorithmus basiert auf dem Paper "An Algorithm for Dependency-Preserving Smart
Home Updates" (Zdankin et al, 2020) [1]. Ziel ist es
ist, wie bereits in Kapitel 1.1 erwähnt, Abhängigkeiten zwischen Geräten zu untersuchen und basierend auf diesen Abhängigkeiten
Updatekonfigurationen zu ermitteln, welchen den größtmöglichen Nutzen erbringen. 
\begin{enumerate}
\item Im ersten Schritt wird eine Übersicht über die aktuelle Updatekonfiguration erstellt. Eine Updatekonfiguration ist eine Zusammenfassung
aller Geräte und ihrer aktuellen Firmwareversion. Dadurch erhält man die Information, welche Dienstleistungen die Geräte aktuell anbieten.
Zusätzlich wird eine Liste aller möglichen Updates eines Geräts erstellt. 
\item Der zweite Schritt ist eine Optimierung bei der versucht wird die Anzahl der potenziellen Updates zu verringern. Dazu werden dominierte 
Updates herausgefiltert und im weiteren Verlauf nicht mehr in Betracht gezogen. Wie genau dies geschieht und welche Updates dominiert sind, wird
in Abschnitt 3.2.2 erläutert.
\item Im dritten Schritt wird mit Hilfe der Übersicht aller Geräte und Updates ein Updatekonfigurationsgraph erstellt. Dieser Graph kann in der
Theorie als nichtdeterministischer endlicher Automat realisiert werden, bei dem jeweils ein Zustand eine Updatekonfiguration darstellt. In Abbildung 3.1
ist ein Ausschnitt eines Graphen von drei Geräten mit jeweils zwei Updates zu sehen. Der Start-Zustand zeigt, dass sich Gerät 1 auf Version 0, Gerät 2 auf Version
2 und Gerät 3 auf Version 4 befindet. Die anderen Zustände bilden die anderen Möglichkeiten auf die man die Geräte updaten kann, wodurch andere
Updatekonfigurationen entstehen.

\begin{figure}[h]
\centering
\includegraphics{"images/UCG"}
\caption{Ausschnitt eines Updatekonfigurationgraphs, (Zdankin et al, 2020, An Algorithm for Dependency-Preserving Smart Home Updates) }
\label{fig:Prob1:MEA}
\end{figure}
\item Im vierten Schritt werden alle Updatekonfigurationen dahingehend überprüft, ob sie eine Abhängigkeit verletzen. Ist dies der Fall wird diese
Konfiguration aus dem Updatkonfigurationsgraphen entfernt.
\item Im fünften Schritt werden alle nicht Pareto-optimalen Konfigurationen herausgefiltert. Pareto-optimal wird in Abschnitt 3.2.2 definiert.
\item Im sechsten Schritt werden die Updatekonfigurationen in Subnetzwerke unterteilt. Besteht zum Beispiel keine Abhängigkeit zwischen zwei
Gruppen von Geräten in einem System, kann man diese Gruppen in Subnetzwerke unterteilen. Für diese Gruppen kann man später unabhängig von
anderen Gruppen die passende Updatekonfiguration auswählen.
\item Im letzten Schritt wird jeder Updatekonfiguration ein Rating zugewiesen mit dem Ziel dem Nutzer die Entscheidung für die passende
Updatekonfiguration zu erleichtern.
\end{enumerate}

\section{Generelle Entscheidungen}
Basierend auf den Erkenntnissen aus den Kapiteln 1 und 2 können nun Designentscheidungen bezüglich des Algorithmus getroffen
werden. Ziel ist es den Prototypen eines zukunftsfähigen Tools zu entwickeln, welcher Smart Home Nutzern die Möglichkeit bietet,
Updates vor ihrer Installation zu überprüfen.
Dabei zielt das Design darauf ab unnötige, architekturspezifische Regularien zu vermeiden, um so möglichst simpel
gestaltet werden zu können. Nichtsdestotrotz sollte es möglichst einfach sein, den Algorithmus in verschiedene Smart Home 
Architekturen zu integrieren. Um solch eine einfache Integration zu gewährleisten, wird die Implementation
in Java modular gehalten. Java bietet sich dafür als objektorientierte Programmiersprache an.
Des Weiteren wird der Algorithmus nicht auf der Basis eines realen Smart Home Systems aufgebaut. Dies wäre kontraproduktiv,
da man als Entwickler dazu neigen würde eine systemspezifische Implementation zu entwickeln. Anstatt also ein reales
Smart Home Systems als Grundlage zu nehmen, werden im Folgenden andere Möglichkeiten genauer beleuchtet. 
Eine Möglichkeit ist es, eine Liste mit 
Geräten, aktueller Firmware Version, möglichen Updates und allen dazugehörigen Dienstleistungen zu erstellen. Fügt
man dieser Liste dann noch Abhängigkeiten hinzu, erhält man alle wichtigen Informationen über ein Smart Home System,
die man benötigt. Vorteil einer solchen Liste ist, dass es sehr einfach wäre sie zu erstellen. Für die spätere Evaluation des 
Algorithmus wäre solch eine statische Liste kontraproduktiv, da der Algorithmus somit immer am gleichen 
System getestet werden würde. Dadurch könnten keine validen Aussagen über die Performance des Algorithmus getroffen werden. Aus 
diesem Grund wird ein Generator genutzt, welcher in Abhängigkeit von verschiedenen Parametern Smart Home Systeme kreiert.
So können die Anzahl der Geräte, die Anzahl der möglichen 
Updates pro Gerät, die Anzahl der Dienstleistungen pro Gerät oder auch die Anzahl der Abhängigkeiten zwischen den
Geräten variiert werden. Dies ermöglich den später implementierten Algorithmus an verschiedenen Systemen zu 
evaluieren, um so ein gutes Fazit über Effizienz und Laufzeit treffen zu können. 

\subsection{Generator}
Das Design eines solchen Generators lässt sich in Java folgenderweise gestalten.
Der Generator erzeugt Geräte, welche in Form von Objekten gespeichert werden. Diese Objekte enthalten Informationen über
die Version, die Dienstleistungen und die möglichen Updates des Geräts. Zusätzlich wird noch das Alter der Updates für spätere
Zwecke gespeichert. Idealerweise sollten Updates nach den Spezifikationen des "Semantic Versioning" definiert werden, da es sich
dabei um eine weitverbreitete Art der Versionierung handelt und die Implementation sich damit an aktuelle Standards hält.
Um jedoch im Rahmen dieser Arbeit die Implementierung zu vereinfachen, wird auf diesen Standard verzichtet und Updates werden
im ganzzahligen Bereich von null an hochgezählt. Dementsprechend kommt nach dem Update null das Update eins und so weiter.
\begin{lstlisting}[language=Java]
public class Device {
    int version;
    ArrayList<Integer> services;
    ArrayList<ArrayList<Integer>> updates;
    ArrayList<ArrayList<LocalDate>> updateAge;
}
\end{lstlisting}
Wie bereits in Kapitel 2.1.1 erwähnt unterscheiden sich Smart Home Geräte hinsichtlicher ihrer Definition.
Da ohne klare Definitionen eine Kommunikation zwischen den Geräten nicht
möglich wäre, wird davon ausgegangen, dass alle Geräte und Dienstleistungen unter einem präskriptiven Standard
definiert sind. Die Geräte können also untereinander kommunizieren ohne, dass eine Übersetzung nötig ist, wie es unter einem deskriptiven
Standard der Fall wäre. Rein für die Implementierung des Algorithmus ist es von Vorteil sich an präskriptiven Standards zu orientieren, 
da durch die Homogenität der Geräte und die redundanzfreie Architektur die Komplexität des Algorithmus gemindert werden kann.


Nachdem die Geräte erstellt wurden, müssen Abhängigkeiten zwischen diesen Geräten Abhängigkeiten kreiert werden, um 
ein Smart Home System zu simulieren. Eine Abhängigkeit herrscht zwischen zwei Geräten, wenn ein Gerät Zugriff auf die
Dienstleistungen eines anderen Geräts benötigt. In Abbildung 3.2 sieht man die Abhängigkeit zwischen einer Heizung und einem Thermostat.
Die Heizung greift periodisch auf die Dienstleistung des Thermostats zu, indem es das Thermostat nach der aktuellen Temperatur fragt.
Liegt die Temperatur nun unter einer bestimmten Schwelle würde die Heizung automatisch anspringen und so gewährleisten, dass immer eine
Mindesttemperatur vorhanden ist.

\begin{figure}[h]
\centering
\includegraphics{"images/Heizung"}
\caption{Abhängigkeit zwischen zwei Geräten}
\label{fig:Prob1:MEA}
\end{figure}
\newpage
Die zwei Komponenten Geräte und Abhängigkeiten reichen aus, um ein theoretisches Smart Home System zu bilden, welches für die Anwendung
des Algorithmus völlig ausreicht. Weitere Voraussetzungen sind, dass innerhalb eines solchen Systems keine Updates automatisch
installiert werden und es möglich ist, wie in Kapitel 2.1, beschrieben, sich eine Übersicht über alle Geräte und deren Informationen zu
beschaffen. In realen Systemen müssten diese Informationen in den Metadaten der Geräte
gespeichert werden, um sie leicht abfragen zu können. Da Metadaten bei den meisten
Geräten bereits existieren, sollte es für die Entwickler dieser Geräte kein großes Problem darstellen die Metadaten
um diese Informationen zu ergänzen [1]. Genauere Details, wie zum Beispiel, ob die Kommunikation über eine Cloud läuft, sind nicht
interessant und werden dementsprechend vernachlässigt.

\subsection{Konkretes Design}
Nun steht die Basis für die Implementation des Algorithmus. Im ersten Schritt werden die Updates aller Geräte untersucht. 
Ziel dieses Schrittes ist es, dominierte Updates zu identifizieren und anschließend zu löschen, um so Rechenarbeit einzusparen.
Updates sind dominiert,
wenn ein anderes Update existiert, welches mindestens die gleichen Dienstleistungen anbietet und zusätzlich aktueller ist.

\begin{figure}[h]
\begin{center}
\includegraphics{"images/Versions"}
\caption{Gerät mit 5 Updates (linke Spalte) und den dazugehörigen Dienstleistungen (rechte Spalte) (Zdankin et al, 2020, An Algorithm for Dependency-Preserving Smart Home Updates [1])}
\label{fig:Prob1:MEA}
\end{center}
\end{figure}

\FloatBarrier
In Abbildung 3.3 ist ein Gerät mit fünf verschiedenen Updates/Versionen zu sehen. In diesem Beispiel ist die Version 1.0 von Version 1.1 dominiert,
da Version 1.1 die gleichen Dienstleistungen wie Version 1.0 anbietet und zusätzlich aktueller ist. Das Gleiche gilt für Version 2.0 und 2.1.
Entfernt man die dominierten Updates, bleiben für dieses Gerät drei potenzielle Updates übrig, nämlich Version 1.1, 2.1 und 3.0.
Diese Optimierung an allen Geräten im System durchzuführen spart eine Menge Rechenzeit.
Im nächsten Schritt wird nämlich mit Hilfe des kartesischen Produkts aller restlichen, nicht-dominierten Updates ein Updatekonfigurationsgraph
erstellt. Dieser Graph beinhaltet alle möglichen Updatekonfigurationen, die im System installiert werden könnten. Seien also in einem System zehn
Geräte mit je zwei Updates bedeutet dies, dass es \( 2^{(10)} = 1024\) verschiedene Updatekonfigurationen geben würde. Bei drei Updates pro Gerät
wären es schon \( 3^{(10)} = 59049\) verschiedene Konfigurationen. Dies macht deutlich, dass jedes im ersten Schritt als dominiert erkanntes
Update die Größe des Updatekonfigurationsgraphs erheblich verringern kann. Wie bereits in Abschnitt 3.1 erwähnt, kann der Updatekonfigurationsgraph in der Theorie als
nichtdeterministischer endlicher Automat realisiert werden. In Java wird dieser Automat in Form einer mehrdimensionalen Arrayliste dargestellt.
In Abbildung 3.4 sind zwei Geräte eines Smart Home Systems zu sehen, eine Heizung und ein Thermostat. 
Die Heizung ist momentan auf Version 0 mit den Services  \{1,2\} und besitzt keine Updates.
Das Thermostat ist momentan auf Version 0 mit den Services  \{5,7\} und kann auf Version 1 mit den Services \{6,8\} geupdatet werden.
Der Updatekonfigurationsgraph sieht in Java dann folgendermaßen aus:
\begin{figure}[!h]
\begin{center}
\includegraphics{"images/UCG Java"}
\caption{Updatekonfigurationsgraph in Form einer List in Java}
\label{fig:Prob1:MEA}
\end{center}
\end{figure}
Nun wird der Updatekonfigurationsgraph nach ungültigen Konfigurationen durchsucht. Eine Konfiguration ist ungültig, wenn
Dienstleistungen fehlen, die eigentlich aufgrund einer Abhängigkeit nötig wären. 
Schauen wir uns wieder Abbildung 3.4 an und legen dabei
fest, dass die Heizung Zugriff auf die Dienstleistung mit der Nummer 7 vom Thermostat benötigt. Dies würde bedeuten, dass die Konfiguration
\{\{1,2\},\{5,7\}\} gültig und die Konfiguration \{\{1,2\},\{6,8\}\} ungültig ist. Alle ungültigen Konfigurationen können ignoriert werden, wodurch
die Performance des Algorithmus wieder gesteigert werden kann.
Nachdem bereits dominierte Updates und ungültige Updatekonfigurationen entfernt wurden, werden im nächsten Schritt dominierte Updatekonfigurationen
herausgefiltert, um so wieder die
Performance des Algorithmus zu steigern. Eine Updatekonfiguration ist dominiert, wenn eine weitere Konfiguration existiert, in der alle Updates entweder
gleich alt oder neuer sind. In Abbildung 3.5 ist der Ausschnitt eines Updatekonfigurationsgraphen von drei Geräten mit jeweils drei Updates zu sehen, wobei die Updates
nach ihrer Aktualität von links nach rechts sortiert sind. Bei Gerät1 also ist das Update \{\{1,3\} aktueller als Update \{\{1,2\} und das Update \{\{1,4\} ist aktueller
als Update \{\{1,3\}. Die drei Zustände stehen für je eine Updatekonfiguration, die Pfeile sollen andeuten, dass es möglich ist von einer Updatekonfiguration auf
eine andere upzudaten. Bei genauerer Betrachtung der einzelnen Updatekonfigurationen ist zu erkennen, dass die grüne Konfiguration von der blauen dominiert wird,
von der roten jedoch nicht. In der blauen Konfiguration besitzen nämlich alle Geräte eine aktuellere Version, in der roten Konfiguration hingegen ist die Version des
zweiten Geräts eine Ältere (\{\{6,7\} ist älter als \{\{6,8\}). Des Weiteren ist zu erkennen, dass die rote Konfiguration ebenfalls von der blauen dominiert
wird, denn auch hier besitzen alle Geräte eine aktuellere Version.
\begin{figure}[!h]
\begin{center}
\includegraphics{"images/Dominierte Konfig"}
\caption{Dominiere Updatekonfigurationen}
\label{fig:Prob1:MEA}
\end{center}
\end{figure}

Durch das Filtern dominierter Updatekonfiguration wird die Anzahl potenzieller Konfigurationen signifikant verringert, wodurch es für den Nutzer einfacher wird, sich
für eine passende Konfiguration zu entscheiden. Um diese Entscheidung noch zusätzlich zu vereinfachen, werden im folgenden Schritt Subnetzwerke erstellt, für die
unabhängig von anderen Subnetzwerken eine Updatekonfiguration ausgewählt werden kann. Es ist nicht unüblich, dass in einem Smart Home System Gruppen vorhanden
sind, die unabhängig von anderen Gruppen kommunizieren. Mögliche Gruppen zwischen denen keine Kommunikation notwendig ist, können zum Beispiel Fernseher+Soundbox
und Backofen+Mikrowelle sein. Hier bietet es sich logischerweise an die möglichen Updatekonfigurationen unabhängig voneinander zu betrachten, da so für
den Nutzer ein einfacherer Überblick verschafft werden kann. 
Ein Subnetzwerk besteht also aus allen Geräten, die durch eine Abhängigkeit
in Verbindung stehen. Diese Verbindung muss keine Direkte sein. Seien zum Beispiel zwei Geräte A und B durch eine Abhängigkeit verbunden und gleichzeitig sei Gerät
A zusätzlich mit einem Gerät C aufgrund einer Abhängigkeit verbunden, dann sind die Geräte B und C ebenfalls auf indirektem Weg miteinander verbunden, obwohl
keine direkte Abhängigkeit zwischen ihnen herrscht. Somit sind alle drei Geräte A, B und C Teil eines Subnetzwerkes.
Der letzte Schritt besteht darin Bewertungen für die übrig gebliebenen Updatekonfigurationen zu erstellen. Zdankin schlägt vor zu diesem Zweck den einzelnen Dienstleistungen ein
Wert zuzuschreiben [1]. Unter einem präskriptiven Standard wäre es theoretisch für Hersteller möglich Daten zu sammeln und so zu ermitteln, welche Dienstleistungen am 
häufigsten verwendet werden. Ein Standard könnte dann zusätzlich zur Definition einer Dienstleistung eine Wertung abgeben, je nach dem wie häufig diese Dienstleistung 
im Gegensatz zu anderen verwendet wird. Nutzerstatistiken allein reichen nicht aus um eine valide Wertung abzugeben, da Dienstleistungen existieren, welche trotz
seltener Nutzung extrem wichtig sein können. Zum Beispiel sollte der Alarm eines Feuermelders nicht aufgrund der seltenen Nutzung dieser Dienstleistung, als weniger wichtig
bewertet werden. Daher bedarf es weiterer Metriken, um eine valide Aussage über den Wert einer Dienstleistung treffen zu können [1].

\section{Implementierter Code}
Die mitgegebene Implementation des beschriebenen Codes besteht aus 4 Klassen:
\begin{itemize}
\item Device.java: Klasse zum Erstellen der Geräte als Objekte.
\item Generator.java: Generator zum erstellen verschiedener Smart Home Systeme
\item SmartHome.java: Klasse in der das Smart Home System gespeichert wird und an der der Algorithmus ausgeführt wird.
\item Algorithm.java: Die Implementierung des Algorithmus wie er beschrieben wurde.
\end{itemize}

Die Klasse Device beinhaltet einen Konstruktor zum Erstellen der Geräte. Der Generator nutzt diesen Konstruktor und erstellt mit der
Funktion \textit{createDevice} Geräte in Form von Objekten. Zwischen den Geräten werden anschließend mit der Funktion \textit{createDependencies} Abhängigkeiten erzeugt,
die in einer Arrayliste  gespeichert werden.
Beim Erstellen der Geräte und der Abhängigkeiten wird auf die Klasse Random aus dem "java.util" package zurückgegriffen, um den Geräten
zufällige Dienstleistungen und Abhängigkeiten zuzuweisen. Die Anzahl der Geräte und der Abhängigkeiten hängt von manuell ausgewählten
Parametern in der Klasse SmartHome ab.
\begin{lstlisting}[language=Java]
public class SmartHome extends GeneratorFinal{
... 
    static int nrOfDevices = 10;
    static int nrOfUpdatesPerDevice = 3;
    static int nrOfServicesPerDevice = 10;
    static int nrOfDependencies = nrOfDevices*2;
...
}
\end{lstlisting}
Die Klasse Algorithmus beinhaltet den eigentlichen Code des implementierten Algorithmus.
Im ersten Schritt werden mit der Methode \textit{isDominated} alle dominierten Updates aus der Arrayliste \textit{SmartHomeDevices} entfernt und
die übrig gebliebenen Updates in die Funktion \textit{cartesianProduct} als Parameter eingefügt, um so den Updatekonfigurationsgraphen zu erstellen.
Danach werden mit der Funktion \textit{deleteBreakingConfigurations} alle ungültigen Konfigurationen aus dem \textit{updateConfigurationGraph} entfernt.
Dafür erstellt die Funktion für jede Konfiguration eine Arrayliste mit allen in dieser Konfiguration vorkommenden Dienstleistungen und überprüft
dann, ob alle notwendigen Dienstleistungen in der Liste vorkommen. Wenn dies nicht der Fall ist, wird die Konfiguration markiert und anschließend
aus dem Updatekonfigurationsgraphen entfernt. 
Zusätzlich existiert eine weitere Arrayliste \textit{updateConfigurationGraphDates}, in der die Alter der Updates zu finden sind. Wird also durch die Funktion eine
Konfiguration in der Liste \textit{updateConfigurationGraph} entfernt, dann werden in der Liste \textit{updateConfigurationGraphDates} die dazugehörigen Alter ebenfalls
gelöscht.
Als Nächstes entfernt die Funktion \textit{paretoOptimal} allle dominierten Updatekonfigurationen, wie in Kapitel 3.2.2 beschrieben.
Die Funktionen \textit{findSubnetworks} und \textit{floodFill} identifizieren anschließend Subnetzwerke basierend auf den Abhängigkeiten
zwischen den Geräten. Nachdem die Subnetzwerke erstellt werden, sollte im nächsten Schritt für jedes Subnetzwerk der passende Updatekonfigurationsgraph erstellt
werden. Dies wurde nicht implementiert, doch theoretisch wäre es möglich, indem die nötigen Informationen aus dem originalen Updatekonfigurationsgraphen extrahiert
werden. Sei ein Updatekonfigurationsgraph für die Geräte 1, 2, 3, 4 und 5 gegeben, so könnte für ein Subnetzwerk, das aus den Geräten 1, 2 und 3 besteht, der Graph kopiert und
alle Informationen über die Geräte 4 und 5 gelöscht werden. Dabei würden im Graph des Subnetzwerkes einige identische Konfigurationen entstehen, die dann noch gelöscht werden
müssten.
Ein Subnetzwerk besteht dann nur noch aus Geräten zwischen denen direkt oder indirekt eine Abhängigkeit besteht.
Im letzten Schritt erstellt die Funktion \textit{rating} für jede Dienstleistung eine Wertung. Die Werte der einzelnen Dienstleistungen werden dann pro
Subnetzwerk summiert, um sie vergleichen zu können. Ein Nutzer könnte sich anschließend für jedes Subnetzwerk für eine Updatekonfiguration entscheiden.




































\chapter{Evaluation}\label{ch:eval}

Im folgenden Kapitel wird die Implementation hinsichtlich ihrer Performance untersucht, wobei der Algorithmus an unterschiedlichen Systemen
getestet wird, die mit Hilfe des Generators erstellt werden. 

\section{Theoretische Performance}
Ein Ziel dieser Evaluation ist es herauszufinden, ob der Algorithmus an aktuellen Smart Home Systemen eine akzeptable Laufzeit vorweist.
Bereits ohne Messungen vorgenommen zu haben, können bei genauerer Betrachtung des Algorithmus Vermutungen zur Komplexität einzelner Schritte aufgestellt werden.
Zum Berechnen des Updatekonfigurationsgraphen muss das kartesische Produkt zwischen allen Geräten und den dazugehörigen Updates gebildet werden. 
Sei ein System mit n Geräten gegeben und pro Gerät existieren s Updates, dann entstehen beim Erstellen des kartesischen Produkts  \( s^n\) mögliche Konfigurationen.
Zur Vereinfachung wird davon ausgegangen, dass alle Geräte in einem System die gleiche Anzahl an Updates besitzen. Die Komplexität dieses Schrittes wäre also
vereinfacht O( \( s^n\)). Für ein "reales" Smart Home System, in dem die Geräte unterschiedlich viele Updates besitzen, läge die Komplexität bei
O($\vert$Gerät1$\vert$*$\vert$Gerät2$\vert$*$\vert$Gerät3$\vert$*...).
Aufgrund dieses exponentiellen Wachstums ist zu vermuten, dass der Rechenaufwand mit jedem weiteren Gerät und jedem Update stark ansteigt.
Bei einem Smart Home System mit zehn Geräten und jeweils fünf Updates pro Gerät entstehen ohne weitere Optimierungen bereits \( 5^{(10)} = 9765625\)
verschiedene Konfigurationen. Sollte der Trend sich fortsetzen, dass die durchschnittliche Anzahl an Geräten pro Smart Home System ansteigt, könnte
es notwendig werden zusätzliche Optimierungsschritte zu entwickeln, um so den Rechenaufwand des Algorithmus zu verringern.
Im Idealfall ergeben die Messungen, dass der Rechenaufwand des Algorithmus so gering ist, dass die Laufzeit auch bei großen Systemen
mit zum Beispiel 30 Geräten vertretbar ist. So kann die Zukunftstauglichkeit des Algorithmus bei immer größer werdenden Smart Home Systemen
gewährleistet werden.
Ein weiterer rechenintensiver Schritt ist das Löschen aller ungültigen Updatekonfigurationen. Zu diesem Zweck muss für jede potenzielle Updatekonfiguration
jeweils eine Liste mit allen angebotenen Dienstleistungen erstellt werden, die dann für weitere Berechnungen des Algorithmus genutzt werden. 
Für die aus dem vorigen Beispiel \( 5^{(10)} = 9765625\) Konfigurationen müssten
also 9765625 Listen erstellt werden, die alle einzeln durchgegangen werden, um so ungültige Updatekonfigurationen zu finden. Auch für diesen Teilschritt 
ist somit eine hoher Rechenaufwand zu erwarten. Während andere
Berechnungsschritte ebenfalls zur Laufzeit des Algorithmus beitragen, ist davon auszugehen, dass sie um einen großen Faktor weniger
rechenintensiv sind und dementsprechend weniger ins Gewicht fallen. Die in diesem Abschnitt aufgestellten Hypothesen werden in den folgenden Abschnitten
überprüft. 

\section{Laufzeitmessungen}
In diesem Kapitel werden mit Hilfe eines Generators verschiedene Smart Home Systeme kreiert, an denen anschließend die Laufzeit des Algorithmus gemessen wird. 
Die Laufzeit des Algorithmus wird stark von der Größe der erstellten Smart Home Systemen abhängig sein. Es stehen folgende Parameter zur Verfügung, um
die Größe eines Smart Home Systems festzulegen: 
\begin{itemize}
\item \#Geräte: Legt die Anzahl der Geräte im System fest
\item \#DienstleistungenProGerät: Legt die Anzahl der Dienstleistungen fest, die ein Gerät anbietet.
\item \#UpdatesProGerät: Legt die Anzahl der Updates pro Gerät fest (aus Implementationsgründen besitzen alle Geräte die gleiche Anzahl an möglichen Updates)
\item \#Abhängigkeiten: Legt die Anzahl der Abhängigkeiten im System fest.
\end{itemize}
Als erstes soll die Frage beantwortet werden, ob der Algorithmus an Systemen mit aktuell durchschnittlicher Größe anwendbar ist.
Zu diesem Zweck wird mit Hilfe des Generators ein durchschnittliches Smart Home System generiert.
Zunächst muss jedoch geklärt werden, wie ein solches Durchschnittssystem aussieht. Laut Statista [19] waren
in amerikanischen Haushalten im Jahr 2020 durchschnittlich zehn Geräte miteinander verbunden. Die durchschnittliche Anzahl an Updates pro Gerät zu
ermitteln, ohne empirische Daten vorliegen zu haben, stellt sich als schwierig dar. Um diesen Parameter nicht willkürlich zu würfeln, wird sich an der Anzahl
der Updates von Apples Iphones orientiert. Seit 2013 erhalten Iphones durchschnittlich zehn Updates pro Jahr[14]. Dabei muss beachtet werden, dass sich die
Veröffentlichung der
Upates über das Jahr verteilen und somit kein Gerät von einen auf den anderen Tag zehn neue Updates erhält. Dennoch ist nicht davon auszugehen,
dass Nutzer täglich überprüfen, ob ein neues Update zur Verfügung steht, so dass sie sich anhäufen. Um täglich manuelle Konfiguration zu vermeiden,
wird davon ausgegangen, dass ein Nutzer regelmäßig alle drei Monate, also insgesamt vier Mal im Jahr, seine Geräte updatet. So sollte ein Gerät durchschnittlich
zwischen zwei und drei (10 Updates pro Jahr / 4 Mal updaten pro Jahr = 2.5)  möglichen Updates besitzen. 
Nun stellt sich noch die Frage wie viele Dienstleistungen
ein Gerät durchschnittlich anbietet und wie viele Abhängigkeiten zwischen den Geräten herrschen. Auch dies ist ohne empirische Daten schwierig zu beurteilen.
Bei Smart Home Geräten handelt es sich meist um simple Geräte, die nur wenigen Zwecken dienen. So kann zum Beispiel ein Thermostat nur die Temperatur messen,
eine Heizung heizen, eine Sicherheitskamera aufnehmen und vielleicht einen Alarm auslösen. Die meisten Geräte bieten also nur wenige Dienstleistungen. Diese
Annahme beruht jedoch auf eigenen Erfahrungen und ist somit nicht wissenschaftlich belegt. Daher wird für das Erstellen der Geräte davon ausgegangen, dass sie durchschnittlich
fünf Dienstleistungen anbieten, um so die Größe eines Systems nicht zu unterschätzen. Aufgrund der geringen Anzahl an Dienstleistungen pro Gerät wird davon ausgegangen,
dass die Anzahl der Abhängigkeiten, die ein Gerät besitzt, ebenfalls gering ist. Viele Geräte funktionieren auch vollkommen selbstständig und besitzen keinerlei
Abhängigkeit. Für die Anzahl der Abhängigkeiten wird daher festgelegt, dass ein Gerät durchschnittlich zwei Abhängigkeiten besitzt. Die Einschätzung basiert ebenfalls auf subjektiven
Erfahrungen, weswegen für zukünftige Arbeiten empirischer Daten bedarf.
\begin{figure}[h]
\begin{center}
\includegraphics[width=14cm,height=9cm]{"images/Durchschnitt"}
\caption{Laufzeitmessungen an einem durchschnittlichen System}
\label{fig:Prob1:MEA}
\end{center}
\end{figure}
\FloatBarrier
Anhand der Abbildung 4.1 ist zu erkennen, dass die Berechnungen im Bereich von Millisekunden durchlaufen sind. Die durchschnittliche Laufzeit des Algorithmus liegt
bei 79 Millisekunden mit Ausreißern bei 350 und 21 Millisekunden. In einem Buch zum Thema Usability Engineering [18] hat Jakob Nielsen verschiedene Grenzwerte zur
Reaktionszeit von Webseiten festgelegt. Diese Werte können als Indikator genutzt werden, um die Laufzeit des Algorithmus zu beurteilen. Sie können 
nicht eins zu eins auf die Laufzeit des Algorithmus übertragen werden, da der Algorithmus in einem anderen Kontext als eine Webseite zu betrachten ist.
Während Webseiten täglich von Nutzern besucht werden, wird der Algorithmus nur einige Male im Jahr verwendet werden. Dementsprechend ist davon
auszugehen, dass die Laufzeiten für Nutzer zumutbar sind, solange sie unterhalb den folgenden Grenzen liegen:
\begin{itemize}
\item 0.1 Sekunden: Prozesse, die unter 0.1 Sekunden reagieren, nimmt ein Nutzer ohne Verzögerung wahr.
\item 1.0 Sekunden: Der Nutzer nimmt eine kurze Verzögerung wahr, die meist als nicht störend interpretiert wird.
\item 10 Sekunden: Für bis zu 10 Sekunden wartet ein Nutzer durchschnittlich auf eine Reaktion. Verzögert sich eine Reaktion um mehr als 10 Sekunden wechselt ein Nutzer
zu anderen Dingen bis eine Reaktion erscheint.
\end{itemize}

Des Weiteren fällt auf, dass der Algorithmus bei gleichbleibenden Parametern starke Schwankungen aufweist. Dies liegt daran, dass trotz
gleichbleibender Parameter die Komplexität eines Smart Homes variieren kann. Überträgt man dies auf die reale Welt, erscheint es logisch, dass
zwei Systeme trotz gleicher Größe nicht automatisch gleich komplex sind. In einem System können viele
Updates dominiert sein, während dies im anderen System nicht der Fall ist. Dies führt dazu, dass die Optimierungen an beiden Systemen unterschiedlich erfolgreich
ausfallen, so dass in einem System mehr Berechnungen durchgeführt werden müssen.

\newpage
\subsection{Laufzeit der Teilschritte}
In diesem Abschnitt werden die Laufzeiten der einzelnen Teilschritte untersucht, um so Problemstellen in der Implementation ausfindig zu machen. Abbildung 4.2 zeigt die
Laufzeit der einzelnen Schritte, die in Kapitel 3 beschrieben wurden. Die Messungen wurden wie im vorigen Kapitel an einem durchschnittlichen Smart Home System durchgeführt.
Die genauen Parameter sind der Abbildung zu entnehmen.
Es ist zu erkennen, dass der Schritt "deleteBreakingConigurations" mit 41,3 Millisekunden circa 65 Prozent der Laufzeit ausmacht, während die Laufzeit der anderen Schritte summiert
nur 34 Prozent der gesamten Laufzeit ausmachen. Die Annahme aus Abschnitt 4.1, dass das Löschen der ungültigen Konfigurationen sehr rechenintensiv im Vergleich zu den
anderen Schritten ist, ist somit richtig. Sollte in zukünftigen Arbeiten versucht werden den Algorithmus zu optimieren, dann sollte diesem Schritt besondere Aufmerksamkeit geschenkt
werden.  Das Erstellen des Updatekonfigurationsgraphen (cartesianProduct) nimmt im Vergleich zu den anderen Schritten nicht so viel Zeit in Anspruch wie vermutet, dennoch
fallen die anderen Teilschritte insgesamt, wie erwartet nicht stark ins Gewicht.

\begin{figure}[h]
\begin{center}
\includegraphics[width=14cm,height=10cm]{"images/Teilschritte"}
\caption{Laufzeit einzelner Teilschritte}
\label{fig:Prob1:MEA}
\end{center}
\end{figure}

\subsection{Einfluss einzelner Parameter}
Im weiteren Verlauf werden die Parameter des Smart Home Generators verändert und die damit verbundenen Laufzeitschwankungen gemessen. So kann ermittelt werden,
welchen Einfluss die einzelnen Parameter auf die Laufzeit haben. Der erste Parameter, der unabhängig von den anderen untersucht wird,
ist die Anzahl der Geräte in einem Smart Home System. 
Dafür wird die Entwicklung der Laufzeit bei steigender Anzahl an Geräten im System gemessen. Um den Einfluss der anderen Parameter möglichst gering zu halten,
sind sie, wie in Abbildung 4.3 zu sehen, eingestellt.
In der Abbildung ist zu erkennen, dass mit steigender Anzahl an Geräten die Laufzeit des
Algorithmus ansteigt. Im Bereich von 1-8 Geräten hat das Hinzufügen eines Geräts eine durchschnittliche Laufzeitsteigerung von circa 15 Prozent zur Folge. Ab 9-13 Geräten
ergeben die Messungen beim Hinzufügen eines Geräts eine Laufzeitsteigerung von circa 30 Prozent. Ab 14 Geräten erhöht sich die Laufzeit mit jedem hinzugefügten
Gerät um jeweils circa 60 Prozent. Die Ergebnisse zeigen, dass das Hinzufügen eines Geräts bei kleinen Systemen (unter 13 Geräten) geringe Laufzeiterhöhungen zur Folge haben.
Je größer ein System ist, desto stärker erhöht sich die sich die Laufzeit mit jedem zusätzlich hinzugefügtem Gerät. Die Messungen wurden für bis zu 20 Geräte in einem
System durchgeführt. Aufgrund der Kurve des Graphen ist zu erwarten, dass der Graph bei noch größeren Systemen immer stärker steigen würde.
\begin{figure}[h]
\begin{center}
\includegraphics[width=14cm,height=8.5cm]{"images/Steigende Geräteanzahl"}
\caption{Laufzeit beim Erhöhen der Anzahl der Geräte im System}
\label{fig:Prob1:MEA}
\end{center}
\end{figure}

\FloatBarrier
In Abbildung 4.3 handelt es sich um theoretische Ergebnisse, um den unabhängigen Einfluss des Steigerns eines einzelnen Parameters zu messen. In einem realistischem Smart Home System
ist davon auszugehen, dass mit steigender Anzahl an Geräten auch die Anzahl
der Abhängigkeiten im System steigt. Außerdem liegt die Anzahl der Updates pro Gerät in einem realistischem Smart Home System, wie bereits erwähnt zwischen 2 und 3. Um also
praxisnähere Ergebnisse zu ermitteln, zeigen die Messungen in Abbildung 4.4 die Laufzeit des Algorithmus an realistischeren Systemen.
Es ist zu erkennen, dass das Erhöhen der Anzahl der Geräte einen ähnlichen Effekt hat wie zuvor. Der Unterschied besteht darin, dass
eine starke Laufzeitsteigerung bereits bei weniger Geräten beginnt und noch stärker ansteigt. Aufgrund des starken Anstiegs war es nicht möglich die Anzahl der Geräte wie in Abbildung
4.3 auf bis zu 20 zu steigern. 
Die durchschnittliche Laufzeiterhöhung pro hinzugefügtem liegt bei 1-6 Geräten bei circa 30 Prozent (doppelt so viel wie in den Messungen in Abbildung 4.3).
Ab 7-13 Geräten liegt die Steigerung bei circa 80 Prozent und bei der Erhöhung von 13 auf 14
Geräten hat sich die Laufzeit um 440 Prozent erhöht. Somit kann man zusammenfassen, dass die Laufzeit des Algorithmus bei realistischen Systemen sehr stark ansteigt, wenn
die Anzahl der Geräte in einem System steigt.   
\begin{figure}[h]
\begin{center}
\includegraphics[width=14cm,height=9cm]{"images/Steigende Geräteanzahl real"}
\caption{Laufzeit beim Erhöhen der Anzahl der Geräte in einem realistischem System}
\label{fig:Prob1:MEA}
\end{center}
\end{figure}

\newpage
Nun folgen weitere Messungen an Systemen, bei denen die Anzahl der Updates pro Gerät erhöht werden. Es wird wieder versucht den Einfluss der anderen Parameter auf
die Laufzeit möglichst gering zu halten. Die genauen Parametereinstellungen sind der Abbildung 4.5 zu entnehmen.
Die Abbildung zeigt, dass bei der Verdopplung der Updates pro Gerät die Laufzeit um durchschnittlich 60 Prozent steigt. Runtergerechnet bedeutet dies, dass jedes Update
die Laufzeit um durchschnittlich 6 Prozent erhöht. Es handelt sich um ein exponentielles Wachstum, was im Vergleich zu Abbildung 4.3 jedoch langsamer ansteigt. Die Messungen
zeigen insgesamt, dass das Verdoppeln der Geräte einen stärkeren Einfluss auf die Laufzeit hat, als das Verdoppeln der Updates im System.
Sollte die Komplexität für den Teilschritt "cartesianProduct",wie in Abschnitt 4.1 beschrieben, bei O( \( s^n\)) liegen, könnte dies ein Grund für den stärkeren
Einfluss der Geräte auf die Laufzeit sein.
Beim Bilden des kartesischen Produkts steht die Anzahl der Geräte im Exponent und die Anzahl der Updates in der Basis. Dementsprechend lässt das Hinzufügen eines
Geräts den Updatekonfigurationsgraphen schneller wachsen, als das Hinzufügen eines Updates. Dies klingt zunächst widersprüchlich zu der Erkenntnis aus Abbildung 4.2, dass
der Teilschritt "cartesianPrdouct" nur wenig Laufzeit benötigt. Die Schritte "cartesianProduct" und "deleteBreakingConfigurations" müssen jedoch zusammenhängend betrachtet
werden, da das Ergebnis des kartesischen Produkt als Eingabe für die Funktion "deleteBreakingConfigurations" genutzt wird.
\FloatBarrier
\begin{figure}[h]
\begin{center}
\includegraphics[width=14cm,height=9cm]{"images/Steigende Updatezahl"}
\caption{Laufzeit beim Erhöhen der Anzahl der Updates pro Gerät im System}
\label{fig:Prob1:MEA}
\end{center}
\end{figure}
\newpage
In Abbildung 4.5 wurden die Messungen wieder an einem unrealistischen System vorge-
nommen, um den Einfluss des Steigerns eines einzelnen Parameters zu messen. 
Also werden im Folgenden Messungen an einem realen System durchgeführt. Aufgrund der Implementation des Generators besitzen alle Geräte die gleiche Anzahl
an Updates. Dementsprechend ist es nicht möglich einem System mit zehn Geräten ein einzelnes Update hinzuzufügen. Für zukünftige Messungen wäre es von Vorteil
den Generator entsprechend zu ändern, so dass dies möglich ist. Erhöht man also den Parameter \#UpdatesProGerät in einem System mit zehn Geräten um eins, bedeutet
dies, dass dem System insgesamt zehn Updates hinzugefügt werden. Bei einem System mit 5 Geräten sind es dementsprechend fünf Updates.
Dementsprechend ist Abbildung 4.6 zu sehen, dass die Laufzeit an einem System mit 10 Geräten schneller ansteigt, als in
einem System mit 5 Geräten. Nichtsdestotrotz ist bei beiden Systemen ein exponentielles Wachstum zu erkennen, welches nicht wie in Abbildung 4.5 gleichmäßig steigt.

\begin{figure}[h]
\begin{center}
\includegraphics[width=14cm,height=9cm]{"images/Steigende Updatezahl real"}
\caption{Laufzeit beim Erhöhen der Anzahl der Updates pro Gerät in einem realistischen System}
\label{fig:Prob1:MEA}
\end{center}
\end{figure}

\newpage
\subsection{Einfluss der Parameter im Zusammenspiel}
Als Letztes wird gemessen, welchen Einfluss die beiden Parameter \#Geräte und \#UpdateProGerät auf die Laufzeit haben, wenn man sie gleichzeitig erhöht.
Basierend auf den vorigen Ergebnissen steigt die Laufzeit des Algorithmus, wie zu erwarten mit dem Erhöhen der Parameter an. 
In Abbildung 4.7 zeigen die Messungen, dass die Laufzeit ab 7 Geräten mit je 7 Updates pro Gerät extrem ansteigt. Der Algorithmus weist als auch hier
ein sehr starkes Wachstum ab einer bestimmten Grenze auf.
\begin{figure}[h]
\begin{center}
\includegraphics[width=14cm,height=9cm]{"images/Beides"}
\caption{Laufzeit beim Erhöhen der Anzahl der Geräte und der Anzahl der Updates pro Gerät in einem realistischem System}
\label{fig:Prob1:MEA}
\end{center}
\end{figure}

\newpage
\section{Speicherbedarf}
In Abbildung 4.8 ist der Speicherbedarf des Algorithmus an einem durchschnittlichen System zu sehen. Es ist zu erkennen,
dass der Speicherbedarf im Gegensatz zur Laufzeit keinen starken Schwankungen unterliegt. Er steigt im Vergleich zur
Laufzeit beim Erhöhen der Geräte und dem Erhöhen der Updates pro Gerät nur langsam an. Durchschnittlich liegt er
bei 1.29 MB. Des Weiteren ist in den Abbildungen 4.3 bis 4.7 zu erkennen, dass der Speicherbedarf ein lineares Wachstum
aufweist. Für die Anwendung des Algorithmus bedeutet dies, dass es kein Problem sein sollte die Berechnungen auf einem Smartphone
durchzuführen. Auch bei steigender Größe der Smart Home Systeme sollte der Algorithmus, was den Speicherbedarf angeht, zukunftstauglich
sein.

\begin{figure}[h]
\begin{center}
\includegraphics[width=14cm,height=6cm]{"images/Speicherverbrauch"}
\caption{Speicherbedarfmessungen an einem durchschnittlichen System}
\label{fig:Prob1:MEA}
\end{center}
\end{figure}

\section{PC-Spezifikationen}
Die Laufzeit des Algorithmus ist von vielen Faktoren abhängig. Um die Ergebnisse besser beurteilen zu können, folgt eine Liste der verwendeten
PC Konfiguration:
\begin{itemize}
\item Betriebssystem: Windows 10 Home 
\item CPU: AMD Ryzen 3 3200g (4 Kerne, Taktfrequenz 3.6 GHz) 
\item RAM: 16 GB DDR4 3000 MHz
\item Speichermedium: SANDISK Ultra 3D NVMe SSD (Lesen: 2400 MB/s, Schreiben: 1750 MB/s)
\end{itemize}

Moderne Smartphone besitzen leistungsstarke Prozessoren. Das meist verkaufte Smartphone im Jahr 2021 ist das Apple Iphone 12 mit einem
A14 Apple Bionic Prozessor und 4 GB LPDDR4X (4266 MHz) RAM [15]. Der Prozessor besitzt 2 Hochleistungskerne mit Taktfrequenzen bis zu 3 GHz und
4 weiteren Kernen mit Taktfrequenzen bis zu 1,82 GHz [16]. Die Komponenten des Iphone 12 sind nicht wesentlich leistungsschwacher als die des
verwendeten Computers. Der Algorithmus sollte dementsprechend auf einem modernen Smartphone ebenfalls in angemessener Zeit durchlaufen sein.
Weitere Verbesserungen der Implementation, zum Beispiel Nebenläufigkeit, würden die Performance zusätzlich steigern, so dass der Algorithmus auch
auf leistungsschwächeren Smartphones funktionieren sollte. 



\section{Ergebnisse}
Die Evaluationsschritte zeigen, dass der Algorithmus zum aktuellen Stand der Smart Home Systeme für viele Nutzer ein wichtiges Produkt sein könnte.
Nutzern wäre es möglich mit einer simplen Anwendung auf dem Handy die Updates eines Smart Home System zu steuern.
Die Ergebnisse sind jedoch mit Vorsicht zu betrachten, da die Einstellungen der Parameter auf vielen Annahmen beruhen, für die es noch
keine ausreichende Forschungsgrundlage gibt. Für zukünftige Messungen wäre es von Vorteil empirische Daten über Smart Home Systeme
zu sammeln, um so mit Hilfe des Generators möglichst realitätsnahe Systeme zu kreieren.
Die Evaluation hat auch gezeigt, dass die Zukunftstauglichkeit des Algorithmus aufgrund der schnell steigenden Laufzeit nicht garantiert ist.
Weitere Optimierungen werden vermutlich notwendig sein, um vertretbare Laufzeiten bei immer größer werdenden Smart Home Systemen zu erreichen.
Durch die modulare Implementierung sollte es möglich sein Änderungen am Algorithmus vorzunehmen und ihn so zu verbessern.
Ein potenzieller Ansatz um den Algorithmus zu verbessern, ist die Reihenfolge des Algorithmus zu ändern. So könnte bereits im ersten Schritt nach Subnetzwerken
gesucht werden, auf denen die restlichen Schritte des Algorithmus unabhängig voneinander ausgeführt werden . Dadurch müsste der Algorithmus insgesamt öfter
ausgeführt werden, wodurch der Anschein entsteht, dass die Laufzeit sich insgesamt erhöhen würde. Andererseits könnte es sein, dass die Laufzeit sich verringert,
da durch das Erstellen der Subnetzwerke, die Eingabe für den Algorithmus pro Subnetzwerk kleiner ist, als das gesamte Netzwerk. Die Laufzeitmessungen haben
gezeigt, dass große Eingaben (also große Smart Home Systeme) die Laufzeit des Algorithmus, aufgrund des exponentiellen Anstiegs stark erhöhen. Es besteht also
die Möglichkeit, dass das Ausführen des Algorithmus auf vielen kleinen Systemen effizienter ist, als auf einem großen System.







\chapter{Conclusion}

The conclusion quickly summarizes the results of the paper
in relation to the previously defined goals and hypotheses.
It usually also includes some information on next steps and further research
that might be required or possible on the presented subject.




% Appendix chapters to be put here. They will be enumerated with capital letters
% if you  did not change the \documentclass options.
\appendix
% \include{content/appendix}

%Bibliography
% We strongly recommend to use bibtex to manage your bibliography. It helps you
% structure your references and helps avoiding missing important data for a correct
% quotation. If you have no other idea jabref (http://jabref.sourceforge.net/)
% might be a good idea (Java runtime environment needed).
\printbibliography

% Include the eidesstattliche Versicherung
% spellchecker:disable
%pagenumbering{null}

\

\cleardoublepage

\


\pagestyle{empty}

\selectlanguage{german}
\textbf{Versicherung an Eides Statt}\\

Ich versichere an Eides statt durch meine untenstehende Unterschrift,
\begin{itemize}
\item[-] dass ich die vorliegende Arbeit - mit Ausnahme der Anleitung durch die Betreuer - selbstständig ohne fremde Hilfe angefertigt habe und
\item[-] dass ich alle Stellen, die wörtlich oder annähernd wörtlich aus fremden Quellen entnommen sind, entsprechend als Zitate gekennzeichnet habe und
\item[-] dass ich ausschließlich die angegebenen Quellen (Literatur, Internetseiten, sonstige Hilfsmittel) verwendet habe und
\item[-] dass ich alle entsprechenden Angaben nach bestem Wissen und Gewissen vorgenommen habe, dass sie der Wahrheit entsprechen und dass ich nichts verschwiegen habe.
\end{itemize}
Mir ist bekannt, dass eine falsche Versicherung an Eides Statt nach \S 156 und \S 163 Abs. 1 des Strafgesetzbuches mit Freiheitsstrafe oder Geldstrafe bestraft wird.
\vfill
Duisburg, \today\\
$\overline{\parbox{4.8cm}{(Ort, Datum)}} ~~~~~~~~~~~~~~~~~~~~~~~~~~~ \overline{\parbox{7cm}{(Vorname Nachname)}}$


\end{document}
